\begin{theorem}[共振窗特征多项式判别式在素数 \(2\) 处的强制整除指数]\label{thm:pom-resonance-disc-2adic-thickness-lb}
\leanverified{paper\_pom\_resonance\_disc\_2adic\_thickness\_lb}
对共振窗 \(q=9,\dots,17\)，令 \(P_q(\lambda)\in\ZZ[\lambda]\) 为最小整系数递推的特征多项式（定理 \ref{thm:pom-moment-resonance}），并记 \(d_q:=\deg P_q\)。
则其判别式满足
\[
2^{\,d_q-1}\ \mid\ \Disc(P_q).
\]
在该窗口内 \((d_q)_{q=9}^{17}=(7,9,9,13,11,13,11,13,13)\)，因此 \((d_q-1)_{q=9}^{17}=(6,8,8,12,10,12,10,12,12)\) 给出相应的统一下界指数。
\end{theorem}
\begin{proof}
由注 \ref{rem:pom-charpoly-mod2-shadow} 与表 \ref{tab:fold_collision_charpoly_mod2_shadow_q9_17}，存在整数 \((a_q,b_q)\) 使
\[
P_q(\lambda)\equiv \lambda^{a_q}(\lambda+1)^{b_q}\pmod 2,\qquad a_q+b_q=d_q,
\]
并且 \(a_q\) 为奇数、\(b_q\) 为偶数。
在特征 \(2\) 下求导得
\[
P_q'(\lambda)\equiv a_q\,\lambda^{a_q-1}(\lambda+1)^{b_q}+b_q\,\lambda^{a_q}(\lambda+1)^{b_q-1}
\equiv \lambda^{a_q-1}(\lambda+1)^{b_q}\pmod 2,
\]
其中第二项因 \(b_q\) 为偶数而消失，第一项因 \(a_q\) 为奇数而保留。
从而
\[
\deg \gcd\!\bigl(P_q\bmod 2,\ P_q'\bmod 2\bigr)=(a_q-1)+b_q=d_q-1.
\]
另一方面，对首一整数多项式 \(P_q\) 有 \(\Disc(P_q)=\pm\Res(P_q,P_q')\)。
又设 \(p\) 为素数，若 \(f,g\in\ZZ[\lambda]\) 满足
\(
r=\deg\gcd(f\bmod p,g\bmod p),
\)
则 Sylvester 结果式矩阵在模 \(p\) 下具有至少 \(r\) 维核，故其行列式 \(\Res(f,g)\) 被 \(p^r\) 整除。
将 \(p=2\)、\((f,g)=(P_q,P_q')\) 与 \(r=d_q-1\) 代入即得 \(2^{d_q-1}\mid\Disc(P_q)\)。证毕。
\end{proof}

\begin{theorem}[\texorpdfstring{$q\in\{13,15\}$}{q in {13,15}} 的判别式在素数 \(3\) 处的强制整除指数]\label{thm:pom-resonance-disc-3adic-q13-15}
对 \(q\in\{13,15\}\)，令 \(P_q(\lambda)\in\ZZ[\lambda]\) 如定理 \ref{thm:pom-moment-resonance}。
则成立
\[
3^{8}\ \mid\ \Disc(P_q).
\]
\end{theorem}
\begin{proof}
由注 \ref{rem:pom-charpoly-mod3-selection} 与表 \ref{tab:fold_collision_charpoly_mod3_selection_q13_15}，在 \(\FF_3[\lambda]\) 中有
\[
P_{13}(\lambda)\equiv P_{15}(\lambda)\equiv \lambda^{2}(\lambda-1)^4(\lambda+1)^5\pmod 3.
\]
在特征 \(3\) 下，指数 \(2,4,5\) 均不被 \(3\) 整除，因此对每个线性因子 \((\lambda-\alpha)^e\)（\(\alpha\in\{0,1,-1\}\)）都有
\(
\gcd\bigl((\lambda-\alpha)^e,\ \partial_\lambda(\lambda-\alpha)^e\bigr)=(\lambda-\alpha)^{e-1}.
\)
于是
\[
\gcd(P_q\bmod 3,\ P_q'\bmod 3)=\lambda(\lambda-1)^3(\lambda+1)^4,
\]
其次数为 \(1+3+4=8\)。
同理于定理 \ref{thm:pom-resonance-disc-2adic-thickness-lb} 的结果式整除性推论，有 \(3^{8}\mid\Res(P_q,P_q')\)，而 \(\Disc(P_q)=\pm\Res(P_q,P_q')\)（首一）给出 \(3^8\mid\Disc(P_q)\)。证毕。
\end{proof}

\begin{theorem}[模 \(2\) 的有限比特决定性与最终周期上界]\label{thm:pom-resonance-mod2-finite-bit-automaticity}
\leanverified{paper\_pom\_resonance\_mod2\_finite\_bit\_automaticity}
对共振窗 \(q=9,\dots,17\)，令 \(S_q(m)\) 为相应碰撞矩序列，并记 \(s_q(m):=S_q(m)\bmod 2\in\FF_2\)。
取表 \ref{tab:fold_collision_charpoly_mod2_shadow_q9_17} 给出的 \((a_q,b_q)\)，并令 \(t_q(m):=s_q(m+a_q)\)。
则满足：
\begin{enumerate}
  \item \((E+1)^{b_q}t_q\equiv 0\pmod 2\)，并且存在唯一系数向量 \((c_{q,0},\dots,c_{q,b_q-1})\in\FF_2^{b_q}\) 使得对所有充分大的 \(m\) 有
  \[
  t_q(m)=\sum_{j=0}^{b_q-1}c_{q,j}\binom{m}{j}\quad(\mathrm{mod}\ 2).
  \]
  \item 令 \(k_q:=\lceil\log_2 b_q\rceil\)，则 \(t_q(m)\) 仅依赖于 \(m\bmod 2^{k_q}\)，从而 \(s_q(m)\) 在 \(m\ge a_q\) 上最终周期满足
  \[
  \mathrm{Per}(s_q)\mid 2^{k_q}.
  \]
\end{enumerate}
\end{theorem}
\begin{proof}
由注 \ref{rem:pom-charpoly-mod2-shadow} 已知
\(
E^{a_q}(E+1)^{b_q}S_q\equiv 0\ (\mathrm{mod}\ 2)
\)，从而 \((E+1)^{b_q}t_q\equiv 0\)。
记 \(\Delta:=E+1\)。
则 \(\Delta\binom{m}{j}=\binom{m}{j-1}\)（约定 \(\binom{m}{-1}=0\)），故对 \(j<b_q\) 有 \(\Delta^{b_q}\binom{m}{j}=0\)。
由定理 \ref{thm:pom-mod2-difference-binomial-basis}，\(t_q\) 在充分大 \(m\) 上必可写成 \(\{\binom{m}{j}\}_{0\le j<b_q}\) 的线性组合；其唯一性由这组序列在 \(\FF_2\) 上的线性无关性给出。
最后由 Lucas 定理知，当 \(0\le j<b_q\) 时 \(\binom{m}{j}\bmod 2\) 仅依赖 \(m\) 的低 \(k_q\) 位二进制展开，因此 \(t_q(m)\) 仅依赖 \(m\bmod 2^{k_q}\)，周期整除结论随之成立。证毕。
\end{proof}

\begin{theorem}[\texorpdfstring{$q\in\{13,15\}$}{q in {13,15}} 的模 \(3\) 两相分解与最终周期上界]\label{thm:pom-resonance-mod3-two-phase-period-q13-15}
对 \(q\in\{13,15\}\)，令 \(u_q(m):=S_q(m+2)\bmod 3\in\FF_3\)。
则存在唯一一对多项式 \(A_q(m),B_q(m)\in\FF_3[m]\) 满足
\[
\deg A_q<4,\qquad \deg B_q<5,
\]
并且对所有充分大的 \(m\) 有
\[
u_q(m)=A_q(m)+(-1)^m B_q(m)\quad(\mathrm{mod}\ 3).
\]
因此 \(u_q\) 的最终周期整除 \(18\)，且其偶子序列与奇子序列均为模 \(3\) 的 \(9\)-周期多项式型序列。
\end{theorem}
\begin{proof}
该结论为命题 \ref{prop:pom-moment-mod3-two-phase} 的等价改写；唯一性来自 \(\FF_3[E]\) 中 \((E-1)^4\) 与 \((E+1)^5\) 的互素分解及对应的 Bézout 投影算子，从而分解 \(u_q=u_q^{(+)}+u_q^{(-)}\) 在两主分量上唯一。
周期上界由 Lucas 定理（基 \(3\)）与二项式基展开给出。证毕。
\end{proof}

\begin{theorem}[\texorpdfstring{$q\in\{13,15\}$}{q in {13,15}} 的模 \(6\) 最终周期上界]\label{thm:pom-resonance-mod6-period-q13-15}
\leanverified{paper\_pom\_resonance\_mod6\_period\_q13\_15}
对 \(q\in\{13,15\}\) 定义 \(v_q(m):=S_q(m)\bmod 6\in\ZZ/6\ZZ\)。
则存在函数 \(F_q:\ZZ/72\ZZ\to\ZZ/6\ZZ\) 使得对一切充分大的 \(m\) 有
\[
v_q(m)=F_q(m\bmod 72),
\]
从而模 \(6\) 的最终周期整除 \(72\)。
\end{theorem}
\begin{proof}
由定理 \ref{thm:pom-resonance-mod2-finite-bit-automaticity}，在对齐 \(m\ge a_q\) 后模 \(2\) 的最终周期整除 \(8\)；
由定理 \ref{thm:pom-resonance-mod3-two-phase-period-q13-15}，在对齐 \(m\ge 2\) 后模 \(3\) 的最终周期整除 \(18\)。
取共同起点并用中国剩余分解 \(\ZZ/6\ZZ\cong\ZZ/2\ZZ\times\ZZ/3\ZZ\)，得到模 \(6\) 的最终周期整除 \(\mathrm{lcm}(8,18)=72\)。证毕。
\end{proof}

\begin{theorem}[共振窗特征多项式在 \texorpdfstring{$\ZZ_2$}{Z2} 上的 Hensel 分解与幂零--幺正分裂]\label{thm:pom-resonance-hensel-splitting-z2}
\leanverified{paper\_pom\_resonance\_hensel\_splitting\_z2}
对共振窗 \(q=9,\dots,17\)，令 \(P_q(\lambda)\in\ZZ[\lambda]\) 如定理 \ref{thm:pom-moment-resonance}，并取表 \ref{tab:fold_collision_charpoly_mod2_shadow_q9_17} 的 \((a_q,b_q)\)。
则在 \(\ZZ_2[\lambda]\) 中存在唯一分解
\[
P_q(\lambda)=U_q(\lambda)\,V_q(\lambda),
\qquad \deg U_q=a_q,\ \deg V_q=b_q,
\]
满足
\[
U_q(\lambda)\equiv \lambda^{a_q}\pmod 2,\qquad V_q(\lambda)\equiv (\lambda+1)^{b_q}\pmod 2.
\]
并且存在典范同构
\[
\ZZ_2[\lambda]/(P_q)\ \cong\ \ZZ_2[\lambda]/(U_q)\ \times\ \ZZ_2[\lambda]/(V_q).
\]
在第二因子上令 \(T:=\lambda+1\)，则其模 \(2\) 约化满足 \(T^{b_q}=0\)，从而 \(\lambda=-1+T\) 呈现为一个 \(2\)-进幺正元加上幂零扰动；相应地，由 \(P_q(E)\) 所强制的递推动力学在 \(2\)-进意义下分裂为“有限步幂零部分”与“幺正--幂零差分流”之直和，其模 \(2\) 的长期行为由定理 \ref{thm:pom-resonance-mod2-finite-bit-automaticity} 的二项式基多项式流刻画。
\end{theorem}
\begin{proof}
由注 \ref{rem:pom-charpoly-mod2-shadow}，在 \(\FF_2[\lambda]\) 中有互素分解
\(
P_q(\lambda)\equiv \lambda^{a_q}(\lambda+1)^{b_q}
\)
且 \(\gcd(\lambda^{a_q},(\lambda+1)^{b_q})=1\)。
由 Hensel 引理（互素因子分解可唯一提升）得存在唯一 \(U_q,V_q\in\ZZ_2[\lambda]\) 使得 \(P_q=U_qV_q\) 且满足所示同余与次数条件。
再由中国剩余定理得到商环分裂同构。
对第二因子，令 \(T=\lambda+1\)，则 \(V_q(\lambda)\equiv T^{b_q}\ (\mathrm{mod}\ 2)\) 强制其约化为局部幂零代数，从而 \(\lambda=-1+T\) 给出幺正加幂零的标准形；将该结构转译到伴随算子（或等价的 \(\ZZ_2[\lambda]/(P_q)\)-模）上的移位作用，即得所述动力学分裂与模 \(2\) 的二项式基流描述。证毕。
\end{proof}

\begin{corollary}[模 \(2\) 台阶影子诱导的 \(2\)-进解模规范分裂]\label{cor:pom-resonance-hensel-splitting-sequence-module}
沿用定理 \ref{thm:pom-resonance-hensel-splitting-z2} 的记号，并设
\[
\mathscr M_q^{(2)}
:=
\left\{
f=(f_m)_{m\ge 0}\in \ZZ_2^{\NN}:\ P_q(E)f=0
\right\}
\]
为由 \(P_q(E)\) 湮灭的 \(2\)-进序列模。
取任意 Bézout 关系
\[
A(\lambda)U_q(\lambda)+B(\lambda)V_q(\lambda)=1
\qquad\text{于 }\ZZ_2[\lambda].
\]
则每个 \(f\in\mathscr M_q^{(2)}\) 都存在唯一分解
\[
f=f_{\mathrm{nil}}+f_{\mathrm{unip}},
\]
其中
\[
f_{\mathrm{nil}}:=B(E)V_q(E)f,
\qquad
f_{\mathrm{unip}}:=A(E)U_q(E)f,
\]
并满足
\[
U_q(E)f_{\mathrm{nil}}=0,
\qquad
V_q(E)f_{\mathrm{unip}}=0.
\]
对任意 \(k\ge 1\)，该分裂在模 \(2^k\) 约化后仍保持为直和分解；特别地，模 \(2\) 约化满足
\[
(\,E+1\,)^{b_q}\,\overline{f}_{\mathrm{unip}}=0,
\]
从而 \(\overline{f}_{\mathrm{unip}}\) 落在由 \((E+1)^{b_q}\) 控制的幺正退化塔中。
\end{corollary}
\begin{proof}
由 Bézout 关系得
\[
1=A(\lambda)U_q(\lambda)+B(\lambda)V_q(\lambda).
\]
对任意 \(f\in\mathscr M_q^{(2)}\) 施加移位算子 \(E\) 得
\[
f=A(E)U_q(E)f+B(E)V_q(E)f=f_{\mathrm{unip}}+f_{\mathrm{nil}}.
\]
再由 \(P_q=U_qV_q\) 与多项式交换性，
\[
U_q(E)f_{\mathrm{nil}}
=U_q(E)B(E)V_q(E)f
=B(E)P_q(E)f
=0,
\]
同理
\[
V_q(E)f_{\mathrm{unip}}
=V_q(E)A(E)U_q(E)f
=A(E)P_q(E)f
=0.
\]
若 \(g\) 同时满足 \(U_q(E)g=0\) 与 \(V_q(E)g=0\)，则再用 Bézout 关系得到
\[
g=A(E)U_q(E)g+B(E)V_q(E)g=0,
\]
故上述和分解为直和，唯一性随之成立。

模 \(2^k\) 的分裂只需将定理 \ref{thm:pom-resonance-hensel-splitting-z2} 的商环直积分解逐级约化即可。
最后，由 \(V_q(\lambda)\equiv (\lambda+1)^{b_q}\pmod 2\) 及 \(V_q(E)f_{\mathrm{unip}}=0\) 的模 \(2\) 约化，立得
\(
(\,E+1\,)^{b_q}\overline f_{\mathrm{unip}}=0
\)。
证毕。
\end{proof}

\endinput
